\section{The Iron Ring Ceremony}\label{the-iron-ring-ceremony}

\stanceinfo{\textbf{CFES CESS} [2023-03-02]}{2024-01-02}

\officialstance{The Canadian Federation of Engineering Students believes that the Iron Ring Ceremony is a necessary transitional ceremony for engineering students, but in its current state is outdated and does not fully represent modern age engineering. It is instead limiting the description of the engineering profession and exclusionary in tone, and should be adjusted to affirm all engineers in their Obligation.}

\stanceheading{The Students' Position}
\begin{itemize}
 \item The Iron Ring Ceremony is an essential rite of passage for engineering students, placed in high regard for graduates after the completion of a long and difficult degree.
 \item At its best, the Iron Ring Ceremony can contribute to a sense of belonging and community in engineering.
 \item The original ritual is not inclusive to all engineering students, leaving them disappointed. The ritual restricts the practice of engineering outside of the original forms of work when the ritual was first enacted as well as presents an outdated worldview regarding colonialism, racism and sexism.
 \item Each camp conducts its own ceremony with little oversight from the corporation of seven wardens.
 \item All engineering students should feel included, comfortable and respected during The Ritual of the Calling of the Engineer.
\end{itemize}

\stanceheading{The Issue}

The Iron Ring Ceremony is the culmination of years of hard work for engineering students in Canada, a ritual where the student is called to take an obligation to practice engineering by ethical and professional standards [7]. This causes the ceremony to hold great importance to many undergraduate students. The main function of the ceremony is to make engineering students feel positive and affirmed as obligated engineers.

Currently, some undergraduate students leave the ceremony feeling a lack of accomplishment and belonging in the engineering profession, straying from the original purpose. Apart from student voices claiming discomfort and discontent within the current Ritual of the Calling of the Engineer, both engineers and engineering educators united to create a statement submitted to the Council of the Seven Wardens in 2022 with a movement ``Retool the Ring'' mainly outlining that it fails to embody a comprehensive understanding of engineering ethics, with a lack of clarity and transparency in text, holding outdated and harmful worldviews rooted in colonialism, racism and sexism [8]. They sought for the Council to acknowledge and take responsibility in order to create a new Call of the Engineer that encompasses modern engineering practice.

The largest challenge with the Iron Ring Ceremony is that these issues do not exist everywhere as there are various currently standing forms of the poem, ``The Ritual of the Calling of the Engineer''. Each Camp organizes the Iron Ring Ceremony independently for a certain geographic area of the country. This camp is allowed to run its own version of the ceremony while some have resolved aforementioned issues. Some Camps have opened the ceremony to guests, others have begun to commission new writings to use in the ceremony.

The CFES believes that all engineering students should be called to be an engineer in the ritual ``Calling of the engineer'' and no student should feel disrespected, uncomfortable or diminished due to their gender, age, religion, race or any outside attribute that has no effect on the engineering practice. As a rite of passage, it must be a ritual that enlightens and encourages the engineers of tomorrow, to practice and grow within their engineering profession.

\stanceheading{Recommendations to Partners, Stakeholders, and Other Entities}
\begin{itemize}
 \item The Partners, Stakeholders, and Other Entities are encouraged to collaborate with the CFES to ensure that the recommendations from the membership are considered, in delivery of the Iron Ring Ceremony.
 \item Integrate both current undergraduate and postgraduate student voices into the current panel enacted titled ``Ritual Review Committee'' [9] for the revision of the Iron Ring Ceremony, based on the Corporation of the Seven Wardens.
 \item For the membership to gather direct information from graduates of changes they would like to see either through regional organizations or their own engineering societies.
 \item The CFES calls Engineers Canada to push for information to be given to undergraduate students before they attend the ceremony.
\end{itemize}

\stanceheading{What the CFES is doing}
\begin{itemize}
 \item Working to collaborate with member organizations to share knowledge and understanding about the Iron Ring Ceremony.
 \item Gathering feedback from graduates that have undergone the Iron Ring Ceremony.
\end{itemize}

\stanceheading{What the CFES plans to do}
\begin{itemize}
 \item Working to ensure proper oversight is conducted over the camps from the Corporation of the Seven Wardens, ensuring that all engineering students are able to positively experience the Iron Ring Ceremony.
 \item To survey engineering students across all camps about their experience during the Iron Ring Ceremony.
 \item Workshops/presentations aimed at disseminating information.
 \item Survey for current students and recent grads about the iron ring ceremony.
 \item Compiling pre-existing surveys.
 \item The CFES shall commit to conducting a minimum of five educational presentations or workshops annually, aimed at disseminating comprehensive information about the Iron Ring Ceremony. These sessions will be targeted towards our member base, ensuring widespread understanding of the ceremony.
 \item The Federation will execute a detailed survey initiative, (date TBD) with the objective of gathering insights from a minimum number of participants (TBD), encompassing current students and recent graduates. This survey will specifically probe themes of inclusivity, representation, impact, and logistics, ensuring that the feedback is relevant, broad-based, and actionable. The findings of this survey will be reviewed annually to track progress and evolving student perceptions.
\end{itemize}

\newpage
